%!TEX encoding = UTF8
%!TEX root = 0-notes.tex

\unchapter{Équations}
\fancyhead[C]{\textbf{Chapitre \thechapter~— Équations linéaires}}
\addcontentsline{toc}{section}{Méthodologie}

\dfn{équation}{
	Une \emphindex{équation} en $x$ est une expression de la forme
		%\[ \text{(expression 1 pouvant dépendre de $x$)} = \text{(expression 2 pouvant dépendre de $x$)}. \]
		\[ f(x) = g(x). \]
	Une expression ne dépendant pas de $x$ est dite \emphindex{constante}.
	
	Une équation peut admettre aucune, une, plusieurs, ou une infinité de solutions.
	Ces solutions ne sont pas toujours exprimables avec les opérations usuelles.
}{}

\ex{}{
	L'égalité suivante, valable pour un nombre $x\neq0$ (car $2$ est divisé par $x$),
		\[ x^7 - \frac2x + 12 = 41 \]
	admet deux solutions, $x\approx-0,068965517223732834608$ et $x\approx1,6273802760023052895$.
}{}

\thm{manipulations d'équations}{
	Deux équations sont \emphindex{équivalentes} s'il est possible de passer de l'une à l'autre à l'aide d'une opération inversible.
		%\[ \text{(expression 1)} = \text{(expression 2)} \underbrace{\iff}_{\text{est équivalente à}} \text{(expression 3)} = \text{(expression 4)}. \]
		\[ f(x) = g(x) \underbrace{\iff}_{\text{« est équivalente à »}} \tilde{f}(x) = \tilde{g}(x). \]
	Une solution de la première équation est une solution de la deuxième et vice versa.
	Les deux équations ont le même \emphindex{ensemble de solutions}.
	
	\setlength{\columnseprule}{1pt}
	\begin{multicols}{2}
	Opérations inversibles usuelles :
	\begin{enumerate}
		\item Ajouter un nombre.
		\item Soustraire un nombre.
		\item Multiplier par un nombre non nul.
		\item Diviser par un nombre non nul.
	\end{enumerate}
	\columnbreak
	Opérations non inversible usuelles :
	\begin{enumerate}
		\item Multiplier par 0.
		\item Mettre au carré.
	\end{enumerate}
	\vfill\,
	\end{multicols}
}{}

\nomen{
	Une équation est \emphindex{résoluble pour $x$} si elle est équivalente à une équation du type
		\[ x = \text{(un nombre) ou (un autre nombre) ou $\dots$}. \]
	Ces nombres sont les \emphindex{solutions} de l'équation.
}

\dfn{équation linéaire}{
	Une \emphindex{équation linéaire} est une identité du type
		\[ a  x + b = c  x + d, \]
	où $a, b, c, d \in \R$ sont des constantes réelles, et $x \in \R$ est un nombre réel à trouver à l'aide de l'équation.
}{}

%\mth{
%	Pour résoudre une équation linéaire, on regroupe les multiples de $x$ d'un côté et les constantes de l'autre : 
%		\[ (a-c)  x = d - b. \]
%	Si $a-c$ est non nul, on multiplie les deux côté par son inverse pour déduire 
%		\[ x = (d-b)\times \dfrac{1}{a-c} =  \dfrac{d-b}{a-c}. \]
%	Dans ce cas, l'équation n'a qu'une seule solution et une réponse sous forme de fraction est généralement attendue.
%}

\ex{}{
	Pour résoudre une équation linéaire, il faut 1) ajouter/soustraire pour regrouper les multiples de $x$ d'un côté et les constantes de l'autre (« isoler $x$ ») ; 
	puis 2) multiplier/diviser pour trouver $x$.
	\begin{align*}
		7x - 1 = x + 4
		&&\underbrace{\iff}_{+1}&& 7x = x + 5 
		&&\underbrace{\iff}_{-x}&& 6x = 5 
		&&\underbrace{\iff}_{\times\frac16}&& x = \frac56
	\end{align*}
	Il est aussi possible d'isoler $x$ à droite.
	\begin{align*}
		7x - 1 = x + 4
		&&\underbrace{\iff}_{-4}&& 7x-5 = x 
		&&\underbrace{\iff}_{-7x}&& -5 = -6x 
		&&\underbrace{\iff}_{\times\frac{-1}6}&&\frac{-5}{-6} = x
	\end{align*}
}{}

\newpage
\addcontentsline{toc}{section}{Exercices}
\fancyhead[C]{\textbf{Chapitre \thechapter~— Exercices}}


\exe{, difficulty=1}{
	Pour chaque paire d'équations, ajouter entre elles le symbole
		\begin{enumerate}[itemindent=1.5cm]
			\item[$\implies$] si la première équation implique la seconde ;
			\item[$\impliedby$] si la deuxième équation implique la première ;
			\item[$\iff$] si les deux équations sont équivalentes ;
			\item[$\centernot\iff$] si les deux équations sont indépendantes.
		\end{enumerate}
	Spécifier l'opération appliquée pour obtenir la deuxième équation à partir de la première le cas échéant.
	
	Par abus de notation et pour gagner du temps, l'absence de symbole signifie en général que les équations sont équivalentes, mais il faut toujours faire attention à ce que cela soit bien le cas !
	
	\begin{center}
	\begin{tabular}{ccc|c|c}
		Équation 1 & Symbole & Équation 2 & Opération & \thead{Les équations sont-\\ elles équivalentes ?} \\ \hline
		$3x + 2y = 1$ &  & $6x  + 4y = 2$ & &  \\ \hline
		$-x + 2y = 1$ &  & $-x + 2y - 1 = 0$ & &  \\ \hline
		$3x + 2y = 1$ & & $-3x - 5y = -3$ & &  \\ \hline
		$x^2 = -3x$ &  & $x = -3$ & &  \\ \hline
		$-10x - y = 3$ &  & $0=0$ & &  \\ \hline
		$x-y= 3$ & & $-x + y = -3$ & &  \\ \hline
		$x -y = 0$ & & $x=y$ & &  \\ \hline
		$-x + 2y = 3$ & & $0=0$ & &  \\ \hline
		$2x + 3y + 4 = 2x + 3y + 5$ & & $1=0$ & &  \\ \hline
		$-x - y = 1$ & & $-x - y = 1$ & &  \\ \hline
		$0=0$ & & $23x - 10y = 6$ & &  \\ \hline
		$x + 2y = 3$ & & $x^2 + 2yx = 3x$ & &  \\ \hline
		$x-3y  =0$ & & $-2x+6y = 0$ & &  \\ \hline
		$-2x + 3y = -3$ & & $4x - 6y = -6$ & &  \\ \hline
		$x^2 = x$ & & $x = 1$ & &  \\ \hline
		$x = 3$ & & $x^2 = 9$ & &  \\ \hline
	\end{tabular}
	\end{center}
}{exe:1}{
	\begin{center}
	\begin{tabular}{ccc|c|c}
		Équation 1 & Symbole & Équation 2 & Opération & \thead{Les équations sont-\\ elles équivalentes ?} \\ \hline
		$3x + 2y = 1$ & {$\iff$} & $6x  + 4y = 2$ & {$\times2$} & {Oui} \\ \hline
		$-x + 2y = 1$ & {$\iff$} & $-x + 2y - 1 = 0$ & {$-1$} & {Oui} \\ \hline
		$3x + 2y = 1$ & {$\centernot\iff$} & $-3x - 5y = -3$ & & {Non} \\ \hline
		$x^2 = -3x$ & {$\impliedby$} & $x = -3$ & & {Non} \\ \hline
		$-10x - y = 3$ & {$\implies$} & $0=0$ & {$\times0$} & {Non} \\ \hline
		$x-y= 3$ & {$\iff$} & $-x + y = -3$ & {$\times(-1)$} & {Oui} \\ \hline
		$x -y = 0$ &{$\iff$} & $x=y$ & {$+y$} & {Oui} \\ \hline
		$-x + 2y = 3$ & {$\implies$} & $0=0$ & {$\times0$} & {Non} \\ \hline
		$2x + 3y + 4 = 2x + 3y + 5$ & {$\iff$} & $1=0$ & {$+(-2x-3y)$} & {Oui} \\ \hline
		$-x - y = 1$ & {$\centernot\iff$} & $-x - y = 1$ & & {Non} \\ \hline
		$0=0$ & {$\impliedby$} & $23x - 10y = 6$ & & {Non} \\ \hline
		$x + 2y = 3$ & {$\implies$} & $x^2 + 2yx = 3x$ & {$\times x$} & {Non} \\ \hline
		$x-3y  =0$ & {$\iff$} & $-2x+6y = 0$ & {$\times(-2)$} & {Oui} \\ \hline
		$-2x + 3y = -3$ & {$\centernot\iff$} & $4x - 6y = -6$ & & {Non} \\ \hline
		$x^2 = x$ & {$\impliedby$} & $x = 1$ & & {Non} \\ \hline
		$x = 3$ & {$\implies$} & $x^2 = 9$ & {Mise au carré} & {Non} \\ \hline
	\end{tabular}
	\end{center}
	
}

\filbreak

\exe{difficulty=1}{
	Un élève de Seconde fournit la production écrite suivante.
	Celle-ci présente quelques erreurs.
	
	\fbox{\parbox{\textwidth}{
		\textbf{\oldunderline{Énoncé :}}
		
		Considérons l'équation $x^3 + x^2 + x + 1 = 0$ où $x\in\R$ est un nombre réel encore inconnu.
		\begin{enumerate}
			\item
			Montrer que $(x-1)(x^3 + x^2 + x + 1) = x^4 - 1$.
			\item
			En déduire la ou les solutions de l'équation originelle.
		\end{enumerate}
		
		\textbf{\underline{Réponse de l'élève :}}
		
		\begin{enumerate}
			\item
			Pour développer $(x-1)(x^3 + x^2 + x + 1)$,
			je distribue d'abord le $x$, puis le $-1$, et j'ajoute les résultats.
				\[ x(x^3 + x^2 + x + 1) + (-1)(x^3 + x^2 + x + 1) = x^4 + x^3 + x^2 + x -x^3 -x^2 - x - 1. \]
			Les puissances de $x$ égales se regroupent pour simplifier l'expression
				\[ x^4 + x^3 + x^2 + x -x^3 -x^2 - x - 1 = x^4 - 1, \]
			comme voulu.
			\item
			Un $x\in\R$ pour lequel $x^3 + x^2 + x + 1 = 0$ vérifie donc $0 = x^4 - 1$ et donc
			$x^4  = 1$.
			En appliquant la racine carrée, je trouve $x^2 = 1$ car $x^4 = (x^2)^2$.
			En l'appliquant encore, j'obtiens finalement $x = 1$.
			
			L'unique solution de l'équation est donc $1$.
		\end{enumerate}
	}}
	\begin{enumerate}
		\item 
		Identifier les erreurs commises par l'élève.
		\item
		Répondre à l'énoncé.
	\end{enumerate}
}{exe:erreurs}{
	\begin{enumerate}
		\item
		La première question est bien traitée et ne comporte aucune erreur.	
		\item
		La deuxième question contient deux erreurs : une lors de la résolution de $x^4 = 1$, et l'autre dans sa conclusion.
		
		Premièrement, appliquer la racine carrée est pertinente, mais $\sqrt{x^2} = |x|$, et non pas $x$ !
		La valeur absolue est importante pour ne pas oublier de solution.
		Ainsi, si $x^4 = 1$, alors $|x^2| = 1$. Comme $x^2 \geq 0$ pour tout $x\in\R$, la valeur absolue peut disparaître ici.
		Puis, si $x^2 = 1$, alors $|x| = 1$, et donc $x=1$ ou $x=-1$. Il y a donc deux solutions réelles à l'équation $x^4  = 1$.
		
		Deuxièmement, la conclusion est erronnée car elle présente une erreur de logique.
		En effet, l'impliquation $x^3 + x^2 + x + 1 = 0 \implies x \in \{-1 ; 1\}$ est vraie, mais l'implication inverse ne l'est pas !
		Nous disposons de deux candidats de solutions qu'il faut vérifier. La solution $x=1$ ne fonctionne pas, mais $x=-1$, elle, fonctionne.
		La seule solution est donc $x=-1$.
		
		En multipliant l'équation originelle par $x-1$, nous avons ajouté artificiellement la solution $x=1$.
		Par exemple, les implications suivantes sont vraie
			\[ x - 1 = 0 \implies (x-3)(x-1) = 0 \implies x \in \{3 ; 1\}. \]
	\end{enumerate}
}


\exe{}{
	Trouver le $x\in\Q$ rationnel vérifiant chaque équation suivante, et l'exprimer sous forme de fraction irréductible.

	\begin{multicols}{2}
	\begin{enumerate}[itemsep=10pt]
		\item $3x + 5 = 2x + 7$
		\item $4x - 6 = 3x + 8$
		\item $5x + 10 = 2x + 4$
		\item $-2x + 3 = 4x - 1$
		\item $7x - 5 = 5x + 9$
		\item $6x + 4 = 3x + 12$
		\item $-3x + 2 = -5x - 4$
		\item $8x - 9 = 6x + 3$
		\item $2x + 11 = -3x + 5$
		\item $-4x + 7 = -x + 10$
	\end{enumerate}
	\end{multicols}
}{exe:inter-affine}
{
	\begin{multicols}{2}
	\begin{enumerate}[itemsep=10pt]
		\item $3x + 5 = 2x + 7 \iff x = 2$
		\item $4x - 6 = 3x + 8 \iff x = 14$
		\item $5x + 10 = 2x + 4 \iff x = -2$
		\item $-2x + 3 = 4x - 1 \iff x = \dfrac{2}{3}$
		\item $7x - 5 = 5x + 9 \iff x = 7$
		\item $6x + 4 = 3x + 12 \iff x = \dfrac{8}{3}$
		\item $-3x + 2 = -5x - 4 \iff x = -3$
		\item $8x - 9 = 6x + 3 \iff x = 6$
		\item $2x + 11 = -3x + 5 \iff x = \dfrac{-6}{5}$
		\item $-4x + 7 = -x + 10 \iff x = -1$
	\end{enumerate}
	\end{multicols}
}

\exe{}{
	Trouver le $x \in \Q$ rationnel (tel que chaque dénominateur est non nul) vérifiant chaque équation suivante, et l'exprimer sous forme de fraction irréductible.
	
	\begin{multicols}{2}
	\begin{enumerate}[itemsep=10pt]
		\item $9 = \dfrac3x$
		\item $-4 = \dfrac2x$
		\item $0,2422 = \dfrac1x$
		\item $0,7 = \dfrac{0,1}x$
		\item $3 - \dfrac4x = -9 + \dfrac1x$
		\item $-2 + \dfrac{-2}x = \dfrac3x - 10$
		\item $\dfrac1{x-2} = 3$
		\item $\dfrac3{x+5} = \dfrac7{x+5} - 3$
	\end{enumerate}
	\end{multicols}

}{exe:eq-lin}
{
	\begin{multicols}{2}
	\begin{enumerate}[itemsep=10pt]
		\item $9 = \dfrac3x \iff x = \dfrac13$
		\item $-4 = \dfrac2x \iff x=-\dfrac12$
		\item $0,2422 = \dfrac1x \iff x = \dfrac{10000}{2422} = \dfrac{5000}{1211}$
		\item $0,7 = \dfrac{0,1}x \iff x = \dfrac17$
		\item $3 - \dfrac4x = -9 + \dfrac1x \iff x = \dfrac5{12}$
		\item $\dfrac1{x-2} = 3 \iff x=\dfrac73$
		\item $\dfrac3{x+5} = \dfrac7{x+5} - 3 \iff x=\dfrac{-11}3$
	\end{enumerate}
	\end{multicols}
}


\exe{, difficulty=1}{
	Multipliez l'équation 
		\[ \dfrac{a}b + \dfrac{c}d = x \]
	par $b$ puis $d$ à gauche et à droite et en déduire la règle d'addition des fractions :
		\[ x = \dfrac{ad + cb}{bd}. \]
}{exe:common-denom}{
	En multipliant par $bd$, on trouve
		\[ (bd) x =  bd \left( \dfrac{a}b + \dfrac{c}d \right) = ad + bc, \]
	d'où
		\[ x = \dfrac{ad+bc}{bd}. \]
}


\exe{}{
	Résoudre les équations suivantes pour $x\in\R$ non nul. 
	%Exprimer $x$ sous la forme $q \sqrt{n}$ avec $q\in\Q$ rationnel et $n\in\N$ le plus petit possible.
	\begin{multicols}{4}
	\begin{enumerate}
		\item $\dfrac{x}3 = \dfrac12.$
		\item $ \dfrac7x  =\dfrac12$
		\item $\dfrac6x  =\dfrac{\sqrt{3}}2$
		\item $\dfrac{x}2 =\dfrac1{\sqrt{3}}$
		%\item $\dfrac9x  =\dfrac{\sqrt{6}}5$
		%\item $-\dfrac3x =\sqrt{7}$
		%\item $-\dfrac7{2x}  =\dfrac{\sqrt{11}}3$
		%\item $\dfrac3{5x} =-\sqrt{5}$
	\end{enumerate}
	\end{multicols}
}{exe:aut0}{
	\begin{multicols}{2}
	\begin{enumerate}
		\item 
			\begin{align*}
				\dfrac{x}3 &= \dfrac12 \\
				x &= 3 \cdot \frac12 \\
				x &= \frac32
			\end{align*}
		\item
			\begin{align*}
				\dfrac7x  &=\dfrac12 \\
				7 &= x \cdot \frac12 \\
				x &= 7\cdot2 = 14
			\end{align*}
		\item
			\begin{align*}
				\dfrac6x  &= \dfrac{\sqrt{3}}2 \\
				6 &= x \cdot \dfrac{\sqrt{3}}2 \\
				x &= 6\cdot\dfrac2{\sqrt{3}} \\
				x &= \frac{12}3 \sqrt3 = 4\sqrt3
			\end{align*}
		\item
			\begin{align*}
				\dfrac{x}2 &= \dfrac1{\sqrt{3}} \\
				x &= 2 \cdot \dfrac1{\sqrt{3}} \\
				x &= \frac23 \sqrt3
			\end{align*}
	\end{enumerate}
	\end{multicols}
}

\exe{}{
	Résoudre les équations suivantes pour $x\in\R$ non nul. Exprimer $x$ sous la forme $q \sqrt{r}$ avec $q\in\Q$ rationnel et $r\in\N$ le plus petit possible.
	\begin{multicols}{2}
	\begin{enumerate}
		\item $\dfrac3x = \dfrac73.$
		\item $ \dfrac7x  =\dfrac1x + 2$
		\item $\dfrac{16}x  =\dfrac{\sqrt{3}}2$
		\item $-\dfrac2x =4\sqrt{3}$
		\item $\dfrac9{-x}  =\dfrac{\sqrt{6}}5$
		\item $-\dfrac3x =\dfrac57\sqrt{7}$
		\item $-\dfrac9{2x}  =\dfrac{\sqrt{11}}3$
		\item $\dfrac2{7x} =-\sqrt{14}$
	\end{enumerate}
	\end{multicols}
}{exe:eq-sqrt}{}

\exe{difficulty=1}{
	\begin{multicols}{2}
	La Terre complète son tour autour du Soleil en environ 365,2422 jours. 
	Un calendrier ne peut donc pas contenir 365 jours tous les ans si celui-ci doit être synchronisé avec les saisons.
	
	Si le calendrier contient moins de 365,2422 jours, il sera en avance.
	S'il contient plus de 365,2422 jours, il sera en retard

	\begin{center}
		\includegraphics[page=4]{2-figures.pdf}
	\end{center}
	\end{multicols}
	
	\begin{enumerate}
	\item
	Calendrier julien\footnote{De Jules César (né en 100 av. J.-C., mort en 44 av. J.-C.).}.
	
	En $46$ av. J.-C., Jules César introduit la règle de l'année bissextile.
	\begin{enumerate}
		\item
		Trouver un entier naturel $a \in \N$ non nul tel que
			\[ 365{,}2422 \approx 365 + \dfrac1a.\]

		\item
		En déduire la règle d'année bissextile du calendrier julien.
		
		\item
		Au bout de 1 600 ans, de combien de jours le calendrier est-il en retard ?
	\end{enumerate}
	
	\item
	Calendrier grégorien\footnote{Du pape Grégoire XIII (1502-1585).}.
	
	En 1582, plus de 1 600 ans après l'introduction de la règle d'année bissextile, le calendrier julien prend donc un retard notable.
	Le pape Grégoire XIII promulgue alors une réforme du calendrier, ainsi créant le calendrier grégorien encore utilisé aujourd'hui.
	
	Afin de corriger le problème de dérive du calendrier, une nouvelle règle est ajoutée à la définition d'année bissextile.
	
	\begin{enumerate}
		\item
		Trouver un entier naturel $b \in \N$ tel que
			\[ 365{,}2422 \approx 365 + \dfrac14 - \dfrac{b}{400}.\]

		\item
		En déduire la règle d'année bissextile du calendrier grégorien.
		
		\item
		Au bout de combien d'années le calendrier est-il en retard d'un jour ?
	\end{enumerate}
	
	\item
	Un nouveau calendrier ?
	\begin{enumerate}
		\item
		Trouver un entier naturel $c \in \N$ tel que
			\[ 365{,}2422 \approx 365 + \dfrac14 - \dfrac{3}{400} - \dfrac{c}{4000}.\]

		\item 
		Créer une nouvelle règle d'année bissextile.
		\item
		Calculer le nombre d'années avant que votre calendrier soit en retard d'un jour.
	\end{enumerate}
	\end{enumerate}
}{exe:calendrier}{}


\exe{}{
	On souhaite connaître les solutions de l'équation du deuxième degré d'inconnue $x\in\R$ :
		\begin{align}
			22x^2 - 125x + 22  = 0. \label{eq:1}
		\end{align}
	\begin{enumerate}
		\item Montrer que $22x^2 - 125x + 22 = (2x-11)(11x-2)$.
		\item En déduire l'ensemble des solutions de l'équation \eqref{eq:1}.
	\end{enumerate}
}{exe:eq7}{
	\begin{enumerate}
		\item On part toujours de la forme factorisée qu'on développe par double distributivité :
			\[ (2x - 11)(11x - 2) = 22x^2 - 4x - 121x + 22 = 22x^2 - 125x + 22. \]
		\item On utilise que le produit est nul si et seulement si un des deux facteurs est nul.
		L'ensemble des solutions de \eqref{eq:1} est donc $\left\{ \dfrac{11}2 ; \dfrac2{11} \right\}$.
	\end{enumerate}
}


\exe{}{
	On souhaite connaître les solutions de l'équation du troisième degré d'inconnue $x\in\R$ :
		\begin{align}
			4x^3 - 6x^2 - 2x + 3  = 0. \label{eq:2}
		\end{align}
	\begin{enumerate}
		\item Montrer que $4x^3 - 6x^2 - 2x + 3 = (2x-3) \left(2x^2-1\right)$.
		\item En déduire l'ensemble des solutions de l'équation \eqref{eq:2}.
	\end{enumerate}
}{exe:eq8}{

	\begin{enumerate}
		\item On part toujours de la forme factorisée qu'on développe par double distributivité.
		On utilise aussi que $x^2 \cdot x = x \cdot x \cdot x = x^3$.
			\[ (2x - 3)(2x^2 - 1) = 4x^3 - 2x -6x^2 + 3. \]
		\item On utilise que le produit est nul si et seulement si un des deux facteurs est nul.
		Or l'équation $2x^2 - 1 = 0$ est équivalente à $x^2 = \dfrac12$.
		En prenant la racine et en utilisant que $\sqrt{x^2} = |x|$, on obtient
			\[ |x| = \sqrt{\dfrac12} = \dfrac1{\sqrt2}. \]
		Les deux valeurs $\pm \dfrac1{\sqrt2}$ (notation $\pm$ pour \og plus ou moins \fg) sont donc solutions.
		On en déduit que l'ensemble des solutions de \eqref{eq:2} est $\left\{ \dfrac32 ; \pm \dfrac1{\sqrt2} \right\} = \left\{ \dfrac32 ; \dfrac1{\sqrt2} ; - \dfrac1{\sqrt2} \right\}$
	\end{enumerate}

}


\exe{}{
	On souhaite connaître les solutions de l'équation du quatrième degré d'inconnue $x\in\R$ :
		\begin{align}
			16x^4 - 24x^2 + 9  = 0. \label{eq:3}
		\end{align}
	\begin{enumerate}
		\item Montrer que $16x^4 - 24x^2 + 9 = \left(4x^2-3\right)^2$.
		\item En déduire l'ensemble des solutions de l'équation \eqref{eq:3}.
	\end{enumerate}
}{exe:eq9}{

	\begin{enumerate}
		\item On part toujours de la forme factorisée qu'on développe à l'aide de l'identité remarquable $(a-b)^2 = a^2 + b^2 -2ab$.
		On utilise aussi que $(x^2)^2 = x^2 \cdot x^2 = x\cdot x\cdot x\cdot x = x^4$.
			\[ (4x^2 - 3)^2 = (4x^2)^2 + 3^2 - 2\cdot(4x^2)\cdot3 = 16x^4 + 9 -24x^2. \]
		\item On utilise que le produit est nul si et seulement si un des deux facteurs est nul.
		Ici, les deux facteurs sont les mêmes, car $(4x^2 - 3)^2 = (4x^2 - 3)(4x^2 - 3)$, donc on se remet à résoudre $(4x^2 - 3) = 0$.
		Ceci est équivalent à $x^2 = \dfrac34$, et donc l'ensemble des $x$ vérifiant \eqref{eq:3} est donné par $\left\{ \pm \sqrt{\dfrac34} \right\} = \left\{ \pm \dfrac12\sqrt{3} \right\}= \left\{ \dfrac12\sqrt{3} ;  -\dfrac12\sqrt{3} \right\}$.
	\end{enumerate}

}




\newpage
\addcontentsline{toc}{section}{Solutions}
\fancyhead[C]{\textbf{Chapitre \thechapter~— Solutions}}

\shipoutAnswer



